\documentclass[12pt]{article}
\usepackage[utf8]{inputenc}
\usepackage[T1]{fontenc}
\usepackage{hyperref}
\title{Quiavicuzas Zapotec Language Resources for K-12 Curriculum Development}
\author{Compiled by Claude Sonnet 5.5 (claude-sonnet-5-5) for Andr\'as Kornai}
\date{}
\begin{document}
\maketitle

\section{Introduction}
Quiavicuzas Zapotec (ISO 639-3 zpj; Glottocode quia1235), also called Northeastern Yautepec Zapotec, is an Otomanguean language of Oaxaca, Mexico. It is listed as vigorous, with about 4,000 speakers according to a 1990 census figure. Sources describe it as unwritten. Little material was found.

\section{Primary Educational Resources}
No primers, textbooks or school programmes were found in the search results.

\section{Linguistic and Reference Materials}
\begin{itemize}
\item Wikipedia article: \url{https://en.wikipedia.org/wiki/Quiavicuzas_Zapotec}
\item Wikitongues language page: \url{https://wikitongues.org/languages/zpj/}
\item Joshua Project language page: \url{https://joshuaproject.net/index.php/languages/zpj}
\end{itemize}
No grammar or dictionary was found.

\section{Religious and Community Resources}
Joshua Project lists no Bible or New Testament translation in this language. No community organisation was found.

\section{Digital and Audio Resources}
\begin{itemize}
\item Global Recordings Network, Zapotec, Quiavicuzas (audio Bible lessons and evangelistic messages): \url{https://globalrecordings.net/en/language/zpj}
\end{itemize}

\section{Curriculum Development Context}
The language is reported to be spoken by all in the community, which is a strong base for a native-language curriculum. No written materials were found, so orthography, primers and a lexicon would all have to be developed with the community.
\end{document}